\section{The Extensional Ledger}

Finite Cohen forcing ended with conditions ordered by extension. A condition
carried a finite family of atoms, each with bounded support and a decision
made within the current tolerance. Extension preserved the old commitments and
added sharper ones. Compatibility forbade a later condition from erasing what
an earlier one had already made readable. That was enough to force a continuum
reading from below, but it left a bookkeeping question. When are two
conditions the same condition?

They cannot be the same merely because they were written the same way. A
finite measurement grammar is not allowed to attach truth to presentation.
The order in which atoms were listed, the name assigned to a support, or the
path by which two compatible fragments were merged cannot become new
continuum structure. If two conditions carry the same obligations, then for
the purposes of the ledger they must be identified.

The chapter therefore quotients the finite Cohen ledger extensionally. Two
conditions are extensionally the same when every atom carried by one is
carried by the other, every overlap agrees within the same bound, and every
decision forced by one is forced by the other. The quotient does not say that
all presentations are harmless. It says that presentation counts only when it
changes the finite obligations the condition carries.

In symbols, the section needs no more than the following reminder:
\[
  p \sim q
  \quad\hbox{when}\quad
  p\hbox{ and }q\hbox{ carry the same finite decisions.}
\]
The symbol is not the authority. The authority is the grammar that decides
what can be recovered by refinement. If a distinction between \(p\) and \(q\)
cannot be recovered by any finite atom, support, bound, or decision, then the
distinction is not a measurement distinction. It is quotiented away.

This quotient preserves the finite horizon. An extensional condition is still
finite because it is represented by finitely many carried atoms. It may have
many presentations, but it does not acquire infinitely many active
obligations. The quotient removes duplicate descriptions; it does not add a
hidden continuum behind the descriptions. Finite support remains the rule by
which the ledger stays honest.

Merge also becomes extensional. If two compatible conditions are merged, the
merged condition is read by the obligations carried in the union, not by the
route through which the union was assembled. Merging \(p\) with \(q\) and
then with \(r\) must identify with merging \(q\) with \(p\) and then with
\(r\), provided the same obligations survive. The ledger may remember the
refinement history when the history carries a real decision. It may not treat
mere order of inscription as a new decision.

Extension is read in the same way. A stronger condition extends a weaker one
when it carries every obligation of the weaker condition and adds a further
bounded decision. If a proposed strengthening abandons an old atom on shared
support, it is not an extension. If it merely writes the old atom under a
different name while preserving the same obligation, it is the same old
condition extensionally. The grammar identifies presentation and preserves
commitment.

This is where countability matters. The finite ledger can be counted because
its atoms can be indexed by finite data: support, bound, local value, and the
small list of neighboring obligations against which compatibility is checked.
A finite condition is then a finite list of finitely indexed atoms. Finite lists
of countable data remain countable. Extensional quotienting can only identify
some of those lists with one another. It does not create a larger collection.
The countable skeleton survives.

The Cantor--Godel--Cohen warning illuminates the point if used carefully.
Completions can differ by distinctions that no chosen finite ledger is able
to recover. The present grammar does not deny that such stories can be told.
It denies them the status of forced measurement structure. A completion is
admissible here only where its distinctions can be pressed back down into
finite refinements. Extensionality is the rule that keeps the upward
completion answerable to the downward record.

Precision supplies a second illustration. A spline completion may assign an
intermediate value between anchors. That value is not a new measurement, and
it is not an arbitrary guess. It is a consequence of the admissible completion
carried by the anchors. In the same way, an extensional condition may have a
completed reading, but the completed reading is answerable to the finite
obligations that forced it. If two presentations force the same obligations,
their completed readings are identified.

The extensional ledger is therefore the first lift of the chapter. Finite
conditions become extensional without ceasing to be finite. The quotient
identifies redundant presentations, preserves finite support, and carries
countability through the continuum door. Once the ledger has become
countable in this stricter sense, the next residue can be read. The gauge
commutator will ask not merely what a condition carries, but what order of
transport the count is able to close.
